\chapter{Persyaratan Sistem}

\section{Sistem Operasi}

PakAI dirancang untuk berjalan di lingkungan Unix dan Unix-like. Sistem yang didukung:

\begin{itemize}
  \item \textbf{Linux} (Ubuntu, Debian, Fedora, Arch, NixOS, dan distribusi lainnya)
  \item \textbf{macOS} (Ventura, Sonoma, Sequoia, dan versi lebih baru)
  \item \textbf{WSL2} (Windows Subsystem for Linux 2) di Windows 10/11
\end{itemize}

\section{Perangkat Lunak yang Dibutuhkan}

\subsection{Go (untuk build dari kode sumber)}

Diperlukan Go versi \textbf{1.21 atau lebih baru} jika ingin melakukan kompilasi
sendiri dari kode sumber.

\begin{lstlisting}[style=terminal]
$ go version
go version go1.25.0 linux/amd64
\end{lstlisting}

Unduh Go dari \url{https://go.dev/dl/} atau gunakan manajer versi seperti
\textbf{asdf}, \textbf{mise}, atau \textbf{nix}.

\subsection{Penyedia AI (opsional, sesuai kebutuhan)}

Pasang dan autentikasi penyedia yang ingin dipantau:

\begin{center}
\begin{tabular}{lll}
\toprule
\textbf{Penyedia} & \textbf{Perangkat} & \textbf{Autentikasi} \\
\midrule
Claude & Claude Code CLI & \cmd{claude login} \\
OpenAI Codex & Pi atau Codex CLI & Login di sumber yang dipilih \\
OpenCode Go & OpenCode & Langganan di opencode.ai \\
\bottomrule
\end{tabular}
\end{center}

Tidak perlu memasang semua penyedia. PakAI akan mendeteksi otomatis penyedia
mana yang tersedia.

\section{Persyaratan Jaringan}

Semua penyedia yang didukung memerlukan koneksi internet:

\begin{itemize}
  \item \textbf{Claude} dan \textbf{OpenAI Codex} memerlukan koneksi internet untuk
        mengambil data penggunaan melalui API OAuth.
  \item \textbf{OpenCode Go} memerlukan koneksi internet untuk menjalankan sesi
        OpenCode CLI. PakAI sendiri membaca data penggunaan dari database lokal
        yang dibuat oleh OpenCode, tetapi OpenCode CLI membutuhkan jaringan saat
        digunakan.
\end{itemize}
